% AY2023 制御工学 第1問〔I〕（転記）
% 出典: 03_制御工学/original/03_制御工学/2023/2023se.pdf
% 要約: 台車質量 m（x2）と入力端 A（変位 x1）の間に並列のバネ k・ダンパ d。
%       運動方程式・伝達関数、ステップ応答・最大値、(5)で A を固定し外力 f(t) 印加。

\section*{第1問〔I〕}

右図に示す系に関する次の問い (1)〜(5) に答えよ. なお, 質量を $m$, 粘性係数を $d$,
バネ定数を $k$ とし, $x_1(t)$ は系への入力となる変位, $x_2(t)$ は質量の変位を表す.
また, 初期状態において系は静止しており, 粘性に関しては, ダッシュポットの動作速度に
比例した粘性抵抗が生じるものとする.

\begin{center}
\begin{tikzpicture}[>=Latex]
  % 質量 m（台車, 変位 x2）
  \draw[thick,fill=gray!12] (0,0.6) rectangle (1.2,2.2);
  \node at (0.6,1.4) {$m$};
  \draw[thick,fill=white] (0.3,0.45) circle (0.15);
  \draw[thick,fill=white] (0.9,0.45) circle (0.15);
  \draw[thick] (-0.2,0.3) -- (1.6,0.3);
  \foreach \x in {-0.1,0.1,...,1.5} \draw[thick] (\x,0.3) -- (\x-0.2,0.1);
  % x2(t)
  \draw[->,thick] (0.3,2.45) -- (1.0,2.45);
  \node at (0.55,2.75) {$x_2(t)$};
  % バネ k（上）
  \draw[thick,decorate,decoration={zigzag,segment length=7,amplitude=4}] (1.2,1.6) -- (3.0,1.6);
  \node at (2.1,2.0) {$k$};
  % ダンパ d（下）
  \draw[thick] (1.2,0.95) -- (1.55,0.95);
  \draw[thick] (1.55,0.73) rectangle (2.0,1.17);
  \draw[thick] (1.8,0.95) -- (3.3,0.95);
  \node at (1.95,0.5) {$d$};
  % 入力端（連結棒）と 点 A
  \draw[thick] (3.0,1.6) -- (3.3,1.6) -- (3.3,0.95);
  \draw[thick] (3.3,1.6) -- (3.3,2.2);
  \draw[thick,fill=white] (3.3,0.95) circle (0.06);
  \node at (3.6,0.95) {A};
  % x1(t)
  \draw[->,thick] (3.3,2.2) -- (4.0,2.2);
  \node at (3.55,2.5) {$x_1(t)$};
\end{tikzpicture}
\end{center}

\begin{enumerate}[label=(\arabic*)]
  \item 系の運動方程式を求めたうえで, 入力を $x_1(t)$, 出力を $x_2(t)$ として伝達関数を求めよ.

  \item 各パラメータを $m=1$, $d=3$, $k=2$ とし, 入力として $x_1(t)$ にステップ入力を与えた
        ときの応答 $x_2(t)$ とその定常値を求めよ.

  \item 問い (2) で求めた出力 $x_2(t)$ の最大値とその発生時間を求めよ.

  \item 各パラメータを $m=1$, $d=1$, $k=0$, 初期値をすべて $0$ とし, 入力として $x_1(t)$ に
        ステップ入力を与えたときの出力 $x_2(t)$ とその定常値を求めよ.

  \item 次に, 点 A を固定端に取付け, 各パラメータを $m=1$, $d=0$, $k=1$ としたうえで,
        質量に入力 $f(t)=e^{-t}$（ただし, $t<0$ のとき $f(t)=0$）を与える.
        この条件下において, 初期値を $x_2(0)=1$, $\dot{x}_2(0)=0$ と設定したときの出力
        $x_2(t)$ を求めよ. 加えて, 充分な時間が経過した後の振幅の最大値 $|x_2(t)|$ を求めよ.
\end{enumerate}
